% SOURCE_STATEMENT: "Do sufficiently small perturbations of Kerr initial data for the vacuum Einstein equations have maximal developments with complete future null infinity whose exterior asymptotically approaches another Kerr solution?"
% Preparation and final self-audit: CONTRIBUTING.md prompts C and D; source/audit artifacts retained in the local physics_statements_2026_10_11 batch.
% !TeX program = lualatex
\documentclass[11pt]{article}
\usepackage[margin=25mm]{geometry}
\usepackage{fontspec}
\setmainfont{Latin Modern Roman}
\usepackage{amsmath,amssymb,mathtools}
\usepackage{unicode-math}
\setmathfont{Latin Modern Math}
\usepackage{xurl,hyperref}
\hypersetup{colorlinks=true,urlcolor=blue}
\setcounter{secnumdepth}{0}
\setlength{\parindent}{0pt}
\setlength{\parskip}{5pt}
\setlength{\emergencystretch}{3em}
\newcommand{\sref}[2]{\hyperlink{source:#1:#2}{[S#2]}}
\newcommand{\cataloguescope}[1]{\paragraph{Recorded target.}#1}
\newcommand{\R}{\mathbb R}
\newcommand{\C}{\mathbb C}
\newcommand{\N}{\mathbb N}
\newcommand{\Z}{\mathbb Z}
\newcommand{\Q}{\mathbb Q}
\newcommand{\F}{\mathbb F}
\newcommand{\eps}{\varepsilon}
\begin{document}
\section{Nonlinear Stability of the Full Kerr Family}\label{statement}

\subsection{Definitions and mathematical statement}

\paragraph{Shared notation.}$\mathbb N=\{1,2,\ldots\}$; the standard number fields have their usual meanings. All additional conventions are stated locally.


\paragraph{Definitions and conventions.} An asymptotically flat vacuum initial data set $(\Sigma,\bar g,K)$ satisfies the Einstein constraints and has maximal globally hyperbolic development $(M,g)$ with $\operatorname{Ric}(g)=0$. The Kerr family describes rotating vacuum black holes by mass $m>0$ and specific angular momentum $a$; its subextremal black-hole range is $|a|<m$. Future null infinity $\mathcal I^+$ is the ideal endpoint boundary of outgoing null rays, and $J^-(\mathcal I^+)$ is its causal past. Completeness of $\mathcal I^+$ is the future-completeness condition in the report's asymptotically flat formulation.

\paragraph{Source scope and qualifications.} The primary formulation also requires a nonempty black-hole region and sufficiently fast asymptotic approach in the exterior. Smallness norm and decay topology are intentionally rough in the report. The ordinary subextremal stability target should not silently be enlarged to extremal horizons, where different phenomena occur, nor reduced to small angular momentum alone.

\paragraph{Full source formulation.} The report additionally requires $M\setminus J^-(\mathcal I^+)\ne\varnothing$ and that the metric on the exterior approach a final Kerr metric at a sufficiently fast rate. Those conditions are part of the intended nonlinear stability target.

\paragraph{Catalogue mathematical target.}
\paragraph{Catalogue wording retained for traceability.}
Do sufficiently small perturbations of Kerr initial data for the vacuum Einstein equations have maximal developments with complete future null infinity whose exterior asymptotically approaches another Kerr solution?

\paragraph{Saved literature context.} The supplied catalogue's assessment dated 2026-08-21 records: Nonlinear stability of Kerr has been proved for sufficiently small angular momentum, but the extracted statement concerns the full Kerr family and the general subextremal range remains unresolved. The assessment is retained as dated provenance; it is not a new comprehensive literature or open-status certification.
\subsection{Short English statement}
Show that small perturbations of rotating black holes keep a complete exterior and settle to another Kerr black hole.
\subsection{Sources}
\begin{itemize}
\item[\hypertarget{source:statement:1}{S1}] ULAM.AI and the UnsolvedMath Contributors, UnsolvedMath, catalogue entry. Catalogue curation under CC BY 4.0; definitions and English exposition added locally.
\url{https://huggingface.co/datasets/ulamai/UnsolvedMath}.
\item[\hypertarget{source:statement:2}{S2}] Mathematical Aspects of General Relativity (2009), Oberwolfach Reports, DOI 10.4171/owr/2009/46, pp. 2585--2646; official report source and mathematical context.
\url{https://doi.org/10.4171/owr/2009/46}.
\item[\hypertarget{source:statement:3}{S3}] Official EMS Press PDF of the same Oberwolfach report; PDF page 5, printed page 2589.
\url{https://ems.press/content/serial-article-files/46246}.
\item[\hypertarget{source:statement:4}{S4}] Klainerman--Szeftel, Kerr stability; reference in the saved dated catalogue literature review.
\url{https://arxiv.org/abs/2104.11857}.
\item[\hypertarget{source:statement:5}{S5}] Giorgi--Klainerman--Szeftel, nonlinear estimates; reference in the saved dated catalogue literature review.
\url{https://arxiv.org/abs/2205.14808}.
\end{itemize}
\label{end:statement}
\end{document}
